\section{Diffuse fields and a sharp mean-curvature comparison}
\label{sec:diffuse}

This section begins the analytical proof of stable rigidity. It uses the
classical estimates of Section~\ref{sec:classical-estimates} and can be
read independently of the geometric deduction in Part I. Its aim is to
bound the mean curvature of a zero layer by the defect from a planar
transition, with the strict coefficient needed for finite-density
rigidity. The argument has three steps: an exact stability identity,
estimates uniform after division by a small defect, and a classification
of the positive functions arising in the normalized limit.

Throughout this section the ambient dimension $n\ge2$ is fixed. All constants
are uniform over entire solutions $v\in C^2(\R^n,[-1,1])$ of
\begin{equation}\label{eq:diffuse-equation}
 \Delta v=v^3-v.
\end{equation}
We call such a solution stable when
\[
 Q_v(\zeta):=\int_{\R^n}
   \bigl(|\nabla\zeta|^2+(3v^2-1)\zeta^2\bigr)\,\dd x\ge0
 \qquad\text{for every }\zeta\in C_c^1(\R^n).
\]
The tests are scalar and unrestricted. The local estimates below do not
require stability; we invoke it only for the quadratic-form inequality.
Our normalization is

\[
 W(t)=\frac{(1-t^2)^2}{4},\qquad
 g(s)=\tanh(s/\sqrt2),\qquad G(s)=1-g(s)^2.
\]
Thus $g'=G/\sqrt2$, $g''=g^3-g$, and the two normalization constants used
below are
\begin{equation}\label{eq:diffuse-tension}
 s_0=\int_\R(g')^2\,\dd s=\frac{2\sqrt2}{3},\qquad
 \sigma=\int_\R G^2\,\dd s=2s_0=\frac{4\sqrt2}{3}.
\end{equation}
The strong maximum principle applied to $1-v$ and $1+v$ shows that a
solution attaining either endpoint is identically that endpoint. We henceforth
assume $|v|<1$. Interior elliptic estimates give bounds, depending only on
$n$ and the derivative order, for all derivatives of $v$.

\subsection{The inverse profile and the exact stability identity}

Modica's bound \eqref{eq:diffuse-modica}, proved in
Section~\ref{sec:classical-estimates}, is $|\nabla v|^2\le2W(v)$.
For the variables introduced next it says that the inverse profile is
1-Lipschitz.

Define the inverse-profile variables
\begin{equation}\label{eq:diffuse-variables}
 q=1-v^2,\qquad U=g^{-1}(v),\qquad p=\nabla U,\qquad
 T=\nabla^2U,\qquad D=1-|p|^2.
\end{equation}
Then $0\le D\le1$, $|p|\le1$, and direct substitution gives
\begin{align}
 \Delta U&=-\sqrt2 vD,\label{eq:diffuse-U}\\
 \mathcal L D&=2|T|^2,
 &\mathcal L&=-\Delta+2\sqrt2 vp\cdot\nabla+2q|p|^2,
 \label{eq:diffuse-D}\\
 (-\Delta+3v^2-1)q&=-q^2D.\label{eq:diffuse-q}
\end{align}
In particular, writing $e(v)=|\nabla v|^2/2+W(v)$,
\begin{equation}\label{eq:diffuse-energy}
 W(v)-\frac12|\nabla v|^2=\frac14q^2D,
 \qquad e(v)=\frac14q^2(2-D).
\end{equation}
All positive-order derivatives of $U$ have uniform bounds as well. To see
this without assuming that $v$ stays uniformly away from its endpoints,
on each ball subtract $U$ at its center. The resulting function is uniformly
bounded there because $|p|\le1$. Equation~\eqref{eq:diffuse-U} first gives
uniform interior $W^{2,r}$ bounds for every finite $r$, then $C^{1,\alpha}$
bounds. Its right side is the smooth function
$-\sqrt2 g(U)(1-|\nabla U|^2)$; the derivatives of every translate of $g$
are uniformly bounded. Interior Schauder estimates and induction give the
assertion.

The zeroth-order coefficient in $\mathcal L$ is nonnegative. If $D$ vanishes
at a point, its equation and the strong minimum principle give $D\equiv0$;
then $T\equiv0$. Thus $U(x)=e\cdot x+c$ with $|e|=1$ and $v=g(U)$.
In all subsequent arguments requiring normalization by $D$, we assume
$D>0$ everywhere.

The variables in \eqref{eq:diffuse-variables} put the planar profile
at $D=0$. They also produce a test that detects its failure without
requiring a regular zero set. For a stable solution, use $q\eta$ in
$Q_v$ and integrate by parts with \eqref{eq:diffuse-q}:
\begin{equation}\label{eq:q-stability}
 0\le Q_v(q\eta)
   =\int_{\R^n}\big(q^2|\nabla\eta|^2-q^3D\eta^2\big)\,\dd x.
\end{equation}
Indeed, the cross term and the $\eta^2|\nabla q|^2$ term combine with
the potential into $\int q\eta^2(-\Delta+3v^2-1)q$.
This is an exact identity, including for compactly supported Lipschitz
$\eta$ by approximation. The inverse-profile formulation and this test
also appear in \citet[Lemma~2.1]{LLWWW2026}; the calculation here is
independent of dimension. The constant $v=0$ is unstable: a cutoff
which is one on $B_R$, vanishes outside $B_{2R}$, and has gradient
$O(R^{-1})$ gives $Q_0<0$ for large $R$.

\subsection{Estimates uniform relative to a small defect}

To compare zero layers later, we need estimates that remain useful as
$D$ tends to zero. The first bound below controls bending, the second
controls relative changes of defect, and the exponential rate fixes the
cost of comparing two distant layers. The derivative estimates also
specify what remains uniform when the observation window grows.

\begin{lemma}[Uniform defect estimates]\label{lem:diffuse-local}
There is a constant $C=C(n)$, and for each $\gamma>0$ a constant
$C_\gamma=C(n,\gamma)$, such that every solution with $D>0$ satisfies
\begin{equation}\label{eq:diffuse-local}
 |T(x)|^2\le CD(x),\qquad |\nabla D(x)|\le CD(x),\qquad
 D(y)\le C_\gamma e^{(2\sqrt2+\gamma)|y-z|}D(z)
 \quad(y,z\in\R^n).
\end{equation}
For every $R<\infty$ and integer $k\ge0$ there is $C_{R,k}=C(n,R,k)$ such
that, for every center $z$,
\begin{equation}\label{eq:diffuse-normalized-bounds}
 \left\|\frac{D}{D(z)}\right\|_{C^k(B_R(z))}
 +\left\|\frac{T}{\sqrt{D(z)}}\right\|_{C^k(B_R(z))}
 \le C_{R,k}.
\end{equation}
Moreover $\inf_{B_R(z)}D\ge C_R^{-1}D(z)$.
The dependence on a growing observation radius can be made explicit:
for every $k\ge0$ and $\gamma>0$, with $\kappa=2\sqrt2+\gamma$,
\begin{equation}\label{eq:diffuse-derivative-growth}
 \begin{split}
 |\nabla^kD(x)|&\le C_{k,\gamma}e^{\kappa|x-z|}D(z),\\
 |\nabla^kT(x)|&\le C_{k,\gamma}e^{\kappa|x-z|/2}\sqrt{D(z)}
 \qquad(x,z\in\R^n).
 \end{split}
\end{equation}
\end{lemma}

\begin{proof}
We first prove the pointwise bounds on a fixed ball, then the exponential
comparison, and finally the normalized derivative estimates.

\emph{Fixed-ball control.}
The drift and potential in \eqref{eq:diffuse-D}, including the derivatives
of their coefficients, are uniformly bounded. Apply the weak Harnack
inequality to the nonnegative supersolution $\mathcal LD\ge0$ on a
fixed enclosing ball, for example $B_{20}(z)$. The constant leading
matrix and bounded lower-order coefficients permit an exponent
greater than one; H\"older's inequality therefore gives
\[
 \int_{B_4(z)}D\le C D(z).
\]
The enclosing radius is fixed, so the constant depends only on
dimension. Integrating \eqref{eq:diffuse-D} against a nonnegative
cutoff equal to one on $B_3(z)$ and supported in $B_4(z)$ gives
\begin{equation}\label{eq:diffuse-T-integral}
 \int_{B_3(z)}|T|^2\le C D(z).
\end{equation}
Indeed all derivatives can be moved onto the cutoff and the bounded drift;
in particular $\operatorname{div}(vp)$ is uniformly bounded.

Put $b=2\sqrt2 vp$. Twice differentiating
\eqref{eq:diffuse-U}, and using
$D_i=-2(Tp)_i$ and
$D_{ij}=-2p_k\partial_kT_{ij}-2(T^2)_{ij}$, gives the explicit system
\begin{equation}\label{eq:diffuse-T-system}
 (\Delta-b\cdot\nabla)T_{ij}
 =-qDT_{ij}+\sqrt2 vqD p_i p_j
       -q(p_iD_j+p_jD_i)+2\sqrt2 v(T^2)_{ij}.
\end{equation}
Repeated indices are summed. Its right side is bounded in norm by
$C(D+|T|)$, using $\nabla D=-2Tp$ and the preliminary uniform bound on
$T$. Consequently $F_0=D^2+|T|^2$ satisfies
\[
 (-\Delta+b\cdot\nabla)F_0\le C F_0.
\]
The local mean-value estimate, \eqref{eq:diffuse-T-integral}, and $D^2\le D$
give $|T(z)|^2\le C D(z)$. We may now regard the $D$ equation as a
homogeneous linear equation with bounded potential
$2q|p|^2-2|T|^2/D$. Ordinary Harnack gives uniform local comparability of
$D$. Finally apply interior $W^{2,r}$ estimates, $r>n$, to
\eqref{eq:diffuse-D}: its source $2|T|^2$ is bounded on a fixed ball by
$CD(z)$. Sobolev embedding gives $|\nabla D(z)|\le CD(z)$.
This argument uses only the boundedness of $|T|^2/D$, not a derivative
bound for that quotient.

\emph{The exponential rate.}
Fix $\kappa=2\sqrt2+\gamma$. For
$r=|x-y|>0$,
\[
 \mathcal L(e^{-\kappa r})\le e^{-\kappa r}
 \left[-\kappa^2+\frac{(n-1)\kappa}{r}
       +2\sqrt2\kappa|v|+2(1-v^2)\right].
\]
The maximum of the last two terms, for $|v|\in[0,1]$, is
$2\sqrt2\kappa<\kappa^2$. Hence this is negative for $r\ge R_\gamma$,
where $R_\gamma$ is uniform. Local Harnack, chained a fixed number of
times, provides a uniform $a_\gamma>0$ such that
$a_\gamma D(y)e^{-\kappa r}\le D(x)$ on $r=R_\gamma$.
For any $\epsilon>0$ subtract $\epsilon$ from the left side; this preserves
the subsolution inequality since the potential of $\mathcal L$ is
nonnegative. It is then negative on a sufficiently large outer sphere.
Comparison on annuli and $\epsilon\downarrow0$ show
$D(x)\ge a_\gamma D(y)e^{-\kappa|x-y|}$ outside the inner ball.
Harnack covers the inner ball and interchange of $x,y$ gives the last
bound in \eqref{eq:diffuse-local}.

\emph{Normalized derivatives.}
Fix $z$, put $\epsilon=D(z)$, and set
$d_\epsilon=D/\epsilon$, $t_\epsilon=T/\sqrt\epsilon$. The estimates
already proved bound $d_\epsilon$, $\nabla d_\epsilon$ and
$t_\epsilon$ on each fixed ball. Their equations are
\begin{align*}
 (\Delta-b\cdot\nabla)d_\epsilon
 &=2q|p|^2d_\epsilon-2|t_\epsilon|^2,\\
 (\Delta-b\cdot\nabla)(t_\epsilon)_{ij}
 &=-q\epsilon d_\epsilon(t_\epsilon)_{ij}
   +\sqrt{2\epsilon}\,vq d_\epsilon p_i p_j
   -q\sqrt\epsilon\big(p_i\partial_jd_\epsilon
                         +p_j\partial_id_\epsilon\big)
   +2\sqrt{2\epsilon}\,v(t_\epsilon^2)_{ij}.
\end{align*}
Their coefficients have uniform smooth bounds, and $0<\epsilon\le1$.
First use interior $W^{2,r}$ estimates on both equations and then
Schauder estimates, shrinking the ball at each step. Induction gives
\eqref{eq:diffuse-normalized-bounds}. The lower bound for $D$ follows
from the local Harnack inequality already established.
Finally apply \eqref{eq:diffuse-normalized-bounds} on a unit ball centered
at $x$, so that $|\nabla^kD(x)|\le C_kD(x)$ and
$|\nabla^kT(x)|\le C_k\sqrt{D(x)}$, and use the exponential comparison
between $x$ and $z$. This proves \eqref{eq:diffuse-derivative-growth}; in
particular the constants on logarithmically growing windows need not be
inferred from an unspecified dependence of $C_{R,k}$ on $R$.
\end{proof}

\subsection{Zeros and decay inside either phase}

The weighted identity \eqref{eq:q-stability} will eventually be compared
with integrals over $\Sigma=\{v=0\}$. We therefore need to know that
$\Sigma$ exists, and that the weight $q$ becomes exponentially small
far from it. The following elementary comparison proves both facts;
its open-set version will also apply to individual chambers between
zero layers.

\begin{lemma}[Clearing inside a phase]\label{lem:diffuse-phase}
A solution with $|v|<1$ has a nonempty zero set $\Sigma=\{v=0\}$.
For every $0<\gamma<\sqrt2$ there is $C_\gamma=C(n,\gamma)$ such that
\begin{equation}\label{eq:phase-decay}
 q(x)\le C_\gamma
       e^{-(\sqrt2-\gamma)\dist(x,\Sigma)}.
\end{equation}
More generally the same estimate, with $\dist(x,\partial\Omega)$ in
place of $\dist(x,\Sigma)$, holds for any solution in an open set
$\Omega$ with $0<v\le1$ there, and also with the signs reversed.
\end{lemma}

\begin{proof}
\emph{Clearing on large balls.}
Let $\phi_R$ be the positive first Dirichlet eigenfunction of a ball
$B_R(x)$, normalized by $\phi_R(x)=\max\phi_R=1$. For
$0<c<1$ such that $\lambda_1(B_R)<1-c^2$,
\[
 \Delta(c\phi_R)+c\phi_R-(c\phi_R)^3
 =c\phi_R\big(1-\lambda_1(B_R)-c^2\phi_R^2\big)>0.
\]
If $v>0$ on a neighborhood of the closed ball, small multiples of
$\phi_R$ lie below $v$. Increase the multiple up to $c$. A first
interior contact contradicts the strict displayed inequality at a
minimum of $v-c\phi_R$; boundary contact is excluded by $v>0$ there.
Thus $v(x)\ge c$. Since $\lambda_1(B_R)=R^{-2}\lambda_1(B_1)$,
for each $c<1$ this holds once $R$ is sufficiently large. If an entire
$|v|<1$ solution had no zero, connectedness would give one strict sign;
letting $R\to\infty$ and then $c\uparrow1$ would give $|v(x)|\ge1$,
a contradiction. The same comparison also shows, uniformly,
that $|v|\to1$ when distance to a phase boundary tends to infinity.

\emph{Decay at every rate below $\sqrt2$.}
It suffices to treat a positive phase. Write $w=1-v$;
then $0\le w\le1$ and
\[
 \Delta w=(2-3w+w^2)w.
\]
Fix $k=\sqrt2-\gamma/2$. By the preceding comparison there is a
uniform $R_0$ such that $\Delta w\ge k^2w$ at all points whose distance
to the phase boundary is at least $R_0$. If the distance of $x$ to that
boundary is $d>R_0+1$, this inequality holds throughout $B_{d-R_0}(x)$.
Let
\[
 A_k(r)=\frac1{|\mathbb S^{n-1}|}
             \int_{\mathbb S^{n-1}}e^{kr\theta_1}\,\dd\theta.
\]
The radial function $A_k(|y-x|)/A_k(d-R_0)$ solves
$\Delta a=k^2a$ and equals one on that ball's boundary. Comparison
gives $w(x)\le A_k(d-R_0)^{-1}$. A spherical cap of angular radius
$r^{-1/2}$ about $\theta_1=1$ shows, for $r\ge1$,
$A_k(r)\ge c_k r^{-(n-1)/2}e^{kr}$. Absorbing the polynomial factor
into $e^{\gamma r/2}$, and using $q\le2w$, proves the assertion.
All balls in this proof may be chosen with closure inside $\Omega$ and
then increased by a limit, so no boundary regularity of $\Omega$ is needed.
\end{proof}

\subsection{The limit seen from a small-defect zero}

The preceding estimates allow division by $D$ without losing local
compactness. We now identify the resulting limit and the quantity in it
that measures mean curvature.
At a regular zero orient the normal by $\nu=p/|p|$, and use
$A=\dd\nu|_{T\Sigma}$, $H=\tr A$. This sign convention gives
\begin{equation}\label{eq:diffuse-H}
 H=\frac{\partial_\nu D}{2|p|^2}\qquad\hbox{on }\Sigma_{\rm reg}.
\end{equation}
Indeed $\Delta U=0$ at a zero, while
$\Delta U=T(\nu,\nu)+|p|H$ and
$\partial_\nu D=-2|p|T(\nu,\nu)$.
In particular $|H|\le CD$ and $|A|^2\le CD$ at zeros with $D\le1/2$,
by Lemma~\ref{lem:diffuse-local}.

\begin{lemma}[Small-defect limits]\label{lem:diffuse-compactness}
Let $v_j$ be entire solutions in the fixed dimension, and let $z_j$ be
zeros with $\epsilon_j=D_j(z_j)>0$ and $\epsilon_j\to0$.
Translate and rotate so that $z_j=0$ and $\nu_j(0)=e_n$, and write
$x=(x',s)\in\R^{n-1}\times\R$. After a subsequence,
\begin{equation}\label{eq:diffuse-limits}
 U_j\longrightarrow s,\qquad
 \frac{D_j}{\epsilon_j}\longrightarrow\mathcal D>0,\qquad
 \frac{T_j}{\sqrt{\epsilon_j}}\longrightarrow\mathcal T
 \quad\hbox{in }C^k_{\rm loc}\text{ for every }k\ge0.
\end{equation}
Here $\mathcal D(0)=1$. The tensor $\mathcal T$ has only tangential
components, is independent of $s$, is trace free, and is harmonic in
$x'$. The positive function $F(x',s)=G(s)\mathcal D(x',s)$ obeys
\begin{equation}\label{eq:diffuse-limit-operator}
 \mathcal L_0 F=2G(s)|\mathcal T(x')|^2,
 \qquad \mathcal L_0=-\Delta+2-G(s).
\end{equation}
Moreover
\begin{equation}\label{eq:diffuse-mean-limit}
 \frac{H_j(0)}{\epsilon_j}\longrightarrow\frac12F_s(0).
\end{equation}
\end{lemma}

\begin{proof}
Lemma~\ref{lem:diffuse-local} gives all normalized derivative bounds
on fixed balls and a positive local lower bound for $D_j/\epsilon_j$.
Since $U_j(0)=0$, $p_j(0)=\sqrt{1-\epsilon_j}\,e_n$, and
$T_j\to0$ locally smoothly, the first convergence follows as well.
The identities
\[
 T_jp_j=-\tfrac12\nabla D_j,\qquad
 (p_j\cdot\nabla)T_j=-T_j^2-\tfrac12\nabla^2D_j
\]
give $\mathcal T e_n=0$ and $\partial_s\mathcal T=0$ after division
by $\sqrt{\epsilon_j}$. Equation~\eqref{eq:diffuse-U} gives
$\tr\mathcal T=0$. Twice differentiating it gives
$\Delta T_j=-\sqrt2\nabla^2(v_jD_j)=O(\epsilon_j)$ on every fixed
ball, including every fixed derivative order. Thus $\Delta'\mathcal T=0$.
The limiting $D$ equation is
\[
 (-\Delta+2\sqrt2 g(s)\partial_s+2G(s))\mathcal D
        =2|\mathcal T(x')|^2.
\]
Conjugation by $G$, using $G'/G=-\sqrt2 g$ and
$G''=(3g^2-1)G$, gives \eqref{eq:diffuse-limit-operator}.
Finally \eqref{eq:diffuse-H}, $G(0)=1$, and $G'(0)=0$ give
\eqref{eq:diffuse-mean-limit}.
\end{proof}

\subsection{Separating the even source from the normal derivative}

In the limit just obtained, the source $2G|\mathcal T|^2$ is even
in the normal coordinate. We use that symmetry to remove its
contribution to $F_s(0)$ without discarding its nonnegative mass.
The next lemma treats the inhomogeneous source explicitly. No bound on its
growth at infinity is assumed. The hypothesis that a nonnegative solution
$F$ exists supplies the local bounds needed in the exhaustion.

\begin{lemma}[Positive decomposition and normal derivative]\label{lem:diffuse-positive}
Let $F\ge0$ be a smooth entire function on $\R^{n-1}\times\R$ and suppose
$\mathcal L_0 F=S$, where $S\ge0$ is smooth and
$S(x',s)=S(x',-s)$. There is a decomposition
\begin{equation}\label{eq:diffuse-Riesz}
 F=E+h,\qquad E\ge0,\quad h\ge0,\quad
 \mathcal L_0 E=S,\quad \mathcal L_0 h=0,
 \quad E(x',s)=E(x',-s).
\end{equation}
In particular $E_s(x',0)=0$ and $\nabla' E_s(x',0)=0$.
With integrals permitted to have value $+\infty$,
\begin{equation}\label{eq:diffuse-sharp-positive}
 \frac{|F_s(0)|}{2}
 \le\frac{1}{2\sqrt2}\,\frac1\sigma
                       \int_\R G(s)^2F(0,s)\,\dd s.
\end{equation}
For every $c>1/(2\sqrt2)$ there exists $P=P(c)<\infty$ such that
\begin{equation}\label{eq:diffuse-positive-window}
 \frac{|F_s(0)|}{2}
 \le\frac{c}{\sigma}\int_{-P}^P G(s)^2F(0,s)\,\dd s
\end{equation}
for all such $F$. The choice of $P$ is independent of the source and of $F$.

More precisely, every extreme point of
\[
 \mathcal K=\{h\ge0:\mathcal L_0 h=0\text{ on }\R^n,\ h(0)=1\}
\]
has the form $e^{b\cdot x'}f(s/\sqrt2)$, where
$a=\sqrt{4-2|b|^2}\in[1,2]$. If $a>1$, $f$ is one of
\begin{equation}\label{eq:diffuse-extremals}
 \frac1a e^{ay}(a-\tanh y),\qquad
 \frac1a e^{-ay}(a+\tanh y);
\end{equation}
if $a=1$, $f(y)=\sech y$.
\end{lemma}

\begin{proof}
There are four steps. We remove the source by monotone exhaustion,
classify the extreme positive homogeneous solutions, compute their
normal derivative-to-mass ratios, and then pass to all solutions on a
single finite interval.

\emph{The even inhomogeneous part.}
For each integer $R$ solve $\mathcal L_0 E_R=S$ on $B_R(0)$ with zero
Dirichlet data. The potential $2-G$ lies in $[1,2]$, so comparison gives
$0\le E_R\le F$. Reflection in $s$ preserves the operator, domain and
source, hence $E_R$ is even in $s$. Comparison on nested balls gives
$E_R\le E_{R+1}$ on $B_R$. Interior estimates and the upper bound by $F$
show that their increasing limit $E$ is smooth locally and solves the
stated equation. This proves \eqref{eq:diffuse-Riesz}. This construction
does not interchange an uncontrolled source integral with a limit.

\emph{The homogeneous extreme points.}
It remains to control $h$ in \eqref{eq:diffuse-Riesz}: $E$ contributes
no normal derivative at zero and only increases the weighted mass.
The set $\mathcal K$
is convex and compact in the topology of smooth convergence on compact
sets: Harnack and interior estimates bound every derivative on each ball,
and the normalization is preserved in a limit. Every nonzero nonnegative
homogeneous solution is strictly positive. For a horizontal translation
$t\in\R^{n-1}$, uniform Harnack chains give
\[
 h(x'+t,s)\le C_t h(x',s)\quad\hbox{for all }(x',s),
\]
with $C_t$ independent of the base point. The operator is invariant under
these translations. If $h$ is extreme in $\mathcal K$, choose
$C>C_t$; both $C^{-1}h(x'+t,s)$ and
$h(x',s)-C^{-1}h(x'+t,s)$ are nonnegative homogeneous solutions.
Normalizing these two summands at
the origin writes $h$ as a convex combination in $\mathcal K$.
Extremality therefore makes the translate proportional to $h$.
The positive proportionality factors form a continuous multiplicative
character of $\R^{n-1}$, so they equal $e^{b\cdot t}$ for some real $b$.
Consequently $h(x',s)=e^{b\cdot x'}f(s/\sqrt2)$ and
\begin{equation}\label{eq:diffuse-extremal-ode}
 -f''-2\sech^2 y\,f=-(4-2|b|^2)f.
\end{equation}

A positive solution $f$ gives the ground-state identity
\[
 \int_\R\big(|\zeta'|^2-2\sech^2y\,\zeta^2
                         +(4-2|b|^2)\zeta^2\big)\,\dd y
 =\int_\R f^2\left|\left(\frac\zeta f\right)'\right|^2\,\dd y\ge0
 \quad(\zeta\in C_c^1(\R)).
\]
Use compactly supported approximations to $\sech y$, for which
$(-\partial_y^2-2\sech^2y)\sech y=-\sech y$.
It follows that $4-2|b|^2\ge1$; its upper bound $4$ is immediate.
For $a=\sqrt{4-2|b|^2}>1$, two linearly independent solutions of
\eqref{eq:diffuse-extremal-ode} are
$f_a^\pm(y)=e^{\pm ay}(a\mp\tanh y)$. They are positive. The
coefficient of $f_a^+$ in any positive solution is nonnegative by its
asymptotic behavior at $+\infty$, and the coefficient of $f_a^-$ is
nonnegative by the behavior at $-\infty$. Extremality excludes a
combination with both coefficients positive. At $a=1$, reduction of
order gives the basis
\[
 \sech y,\qquad \sech y\int_0^y\cosh^2t\,\dd t.
\]
The second member dominates with opposite signs at the two ends, so a
positive solution is a multiple of $\sech y$. This proves the stated
necessary classification of extreme points.

\emph{The sharp full-line coefficient.}
We compute the derivative-to-mass ratio for the unnormalized $f_a^+$;
reflection gives the other sign. Define
\[
 M(a)=\frac34\int_\R e^{ay}\sech^4y\,\dd y,\qquad 1\le a\le2.
\]
Integration by parts, with boundary term zero because $a<4$, gives
\[
 \frac1\sigma\int_\R G(s)^2f_a^+(s/\sqrt2)\,\dd s
       =\frac{3a}{4}M(a),\qquad
 \left|\frac{\dd}{\dd s}f_a^+(s/\sqrt2)\Big|_{s=0}\right|
       =\frac{a^2-1}{\sqrt2}.
\]
Thus the ratio of half the derivative to this mass is
\begin{equation}\label{eq:diffuse-ratio}
 \frac{\sqrt2}{3}\frac{a-1/a}{M(a)}.
\end{equation}
The function $M$ is convex. Its value and derivative at $2$ can be
computed without special-function formulas: with $t=(1+\tanh y)/2$,
\[
 M(2)=6\int_0^1t^2\,\dd t=2,\qquad
 M'(2)=3\int_0^1t^2\log\frac{t}{1-t}\,\dd t=\frac32.
\]
For the last equality the two logarithmic integrals are $-1/9$ and
$-11/18$, respectively. The tangent-line inequality for convex $M$
therefore yields
\begin{equation}\label{eq:diffuse-M-bound}
 M(a)\ge\frac32a-1\ge\frac43(a-1/a),\qquad 1\le a\le2,
\end{equation}
because the difference in the second inequality is
$(a-2)(a-4)/(6a)\ge0$. The ratio in \eqref{eq:diffuse-ratio} is at
most $1/(2\sqrt2)$.

\emph{One finite interval for the entire cone of positive solutions.}
The full-line mass need not be continuous in the local smooth topology.
We therefore pass to a finite interval before applying convex compactness. The
normalized functions $f_a^\pm/a$, including $a=1$, form a compact
family, and
\[
 \sech^4y\,\frac{f_a^\pm(y)}{a}\le C e^{-2|y|}
 \quad(1\le a\le2).
\]
Their full weighted masses have a strictly positive uniform lower
bound. Consequently, for every $c>1/(2\sqrt2)$, a single sufficiently
large $P$ makes the two linear inequalities
\[
 \pm\frac12h_s(0)\le\frac c\sigma
                 \int_{-P}^P G(s)^2h(0,s)\,\dd s
\]
true for every extreme point of $\mathcal K$. Both sides are continuous
linear functionals in the local smooth topology. The Krein--Milman
theorem extends these inequalities from the extreme points to their
closed convex hull, which is $\mathcal K$. Scaling also covers every
nonnegative homogeneous solution, including zero. Since $E_s(0)=0$
and $E\ge0$, the same inequality holds for $F=E+h$.
This proves \eqref{eq:diffuse-positive-window}. It implies
\eqref{eq:diffuse-sharp-positive} by bounding the truncated mass by
the full mass and sending $c\downarrow1/(2\sqrt2)$.
\end{proof}

\begin{remark}\label{rem:diffuse-truncation}
The strict relaxation in the finite-window statement is necessary.
For $a=2$ and $b=0$ the full-line ratio equals $1/(2\sqrt2)$, and any
finite truncation strictly decreases its positive denominator. In
particular a finite-window proof cannot retain the exact endpoint
constant merely by invoking local convergence.
\end{remark}

\subsection{Returning to a finite normal segment}

The limit inequality now gives a uniform estimate for the original
solution. For a regular zero $z$ and a finite $P>0$, write
\begin{equation}\label{eq:diffuse-normal-average}
 I_P(z)=\frac1\sigma\int_{-P}^P
       q(z+t\nu(z))^3D(z+t\nu(z))\,\dd t.
\end{equation}
This is a normal average, not the pointwise defect. Its normalization is
chosen so that the limiting mass is $\sigma^{-1}\int G^2F$.
The next lemma proves $|H|\le cI_P$ by first fixing the interval and
only then choosing the defect threshold.

\Needspace{3\baselineskip}
\begin{lemma}[Mean curvature controlled by normal defect mass]
\label{lem:diffuse-mean}
For every $c>1/(2\sqrt2)$ there exist $P=P(c)<\infty$ and
$\delta=\delta(n,c)\in(0,1/2]$ such that every entire solution and every zero
$z$ with $0<D(z)<\delta$ satisfy
\begin{equation}\label{eq:diffuse-mean-window}
 |H(z)|\le\frac c\sigma
     \int_{-P}^P q(z+s\nu(z))^3D(z+s\nu(z))\,\dd s.
\end{equation}
The point is regular and the normal is oriented by $p$.
The inequality also holds with any larger normal interval.
In particular one can fix
\begin{equation}\label{eq:diffuse-strict-mean}
 \frac{1}{2\sqrt2}<c_*<\frac{\sqrt2}{3}
\end{equation}
and obtain constants $P_*,\delta_*>0$ used in the finite-density argument.
\end{lemma}

\begin{proof}
Choose $c_1$ strictly between $1/(2\sqrt2)$ and $c$, and fix a finite
$P\ge1$ supplied by Lemma~\ref{lem:diffuse-positive} for $c_1$.
If no uniform $\delta$ existed for this $P$, there would be solutions
and zeros with $0<\epsilon_j=D_j(z_j)\to0$ violating
\eqref{eq:diffuse-mean-window}. Apply
Lemma~\ref{lem:diffuse-compactness} to translate and rotate them.
On the fixed interval $[-P,P]$, smooth convergence gives
\[
 \frac1{\epsilon_j}\int_{-P}^P q_j(se_n)^3D_j(se_n)\,\dd s
 \longrightarrow\int_{-P}^P G(s)^3\mathcal D(0,s)\,\dd s
 =\int_{-P}^P G(s)^2F(0,s)\,\dd s>0.
\]
The source in \eqref{eq:diffuse-limit-operator} is nonnegative and even
in $s$, so the limiting mean satisfies the same bound with coefficient
$c_1$. The violated inequalities and
\eqref{eq:diffuse-mean-limit} instead give the reverse weak inequality
with coefficient $c$. The displayed strictly positive mass and
$c_1<c$ contradict each other. All compactness constants depend only on
$n$ and the chosen finite interval, so the resulting threshold is
uniform over the entire class of solutions. Finally the integrand is
nonnegative, which proves the assertion about larger intervals.
\end{proof}
